\documentclass[10pt,a4paper]{amsart}
\usepackage[margin=19mm]{geometry}
\usepackage{amsmath,amssymb,mathtools}
\usepackage[scr=rsfs,cal=euler]{mathalfa}
\usepackage{xfrac}
\usepackage[unicode,hidelinks,bookmarks=false]{hyperref}
\newcommand{\ie}{i.e.\ }
\newcommand{\cf}{cf.\ }
\newcommand{\resp}{respectively\ }
\newcommand{\Subsection}{\subsection}
\providecommand{\nobreakdash}{\nobreak\hbox{-}}
\setlength{\emergencystretch}{2em}
\multlinegap0pt
\usepackage{mathtools}
\usepackage[shortlabels]{enumitem}
\newtheorem{lem}{Lemma}
    \DeclareMathOperator{\Id}{Id}
\title{Pushforward dynamics on Wasserstein spaces and measure rigidity}
\author{Douglas Finamore, Andr\'e Magalh\~aes de S\'a Gomes, Christian S. Rodrigues}
\date{Source: arXiv:2609.00451v1 (CC BY 4.0)}
\begin{document}
\maketitle

\noindent\emph{Source and licence.} Pushforward dynamics on Wasserstein spaces and measure rigidity. Version arXiv:2609.00451v1 (CC BY 4.0), distributed by arXiv under the Creative Commons Attribution 4.0 International licence, http://creativecommons.org/licenses/by/4.0/, as recorded in the arXiv licence metadata for this exact version and shown on the abstract page https://arxiv.org/abs/2609.00451v1. Full licence notice: \texttt{licenses/arxiv-2609.00451-CC-BY-4.0.txt}.\par\smallskip

\noindent\textit{Editorial scope.} Counted unit: the article's lem lem:exp.estimates (source0.tex lines 722--730). The article's complete original proof of it is delivered with it (source0.tex lines 2503--2570, one \texttt{proof} environment of 68 lines, which the article places away from the statement). No mathematics is written, repaired or re-derived: the statement and the proof are the article's own lines, in the article's own notation, and no retained line is substituted or re-flowed.  Retained uncounted as the source-local context the proof needs: lines 2497--2499.  Nothing was omitted from any retained span, so the package carries no omission record.  No retained line contains a citation, so the wrapper carries no bibliography and no bibliography entry.  No retained line contains a figure environment, an image-inclusion command or drawing code, so no image is delivered and the package declares no asset dependency.  Editorial formatting only, in addition: an AMS \texttt{amsart} preamble that restates the article's own declarations and macros used by the retained lines, namely the environments \texttt{lem} on separate counters, together with the standard packages those lines need.  The retained environments print in this document as Lemma 1 in order of appearance, whereas the article prints the same items with the numbering of its own full document.  The complete native source of the article as distributed in the arXiv e-print for this exact version is bundled unchanged as \texttt{sources/209029/source0.tex}.
    \begin{equation}\label{eq:triangle.inequality}
        d\left(\exp_x\xi,\exp_x\eta\right)\le\le\lvert\xi\rvert+\lvert\eta\rvert,\qquad \xi,\eta\in T_xM.
    \end{equation}

\begin{lem}
\label{lem:exp.estimates}
    There are constants \( \Lambda\ge1 \) and \( C_\ast\ge0 \), both depending only on \( M \), such that for every \( x\in M \), the exponential \(\exp_x: B(0, r_1) \to M\) is \(\Lambda\)-Lipschitz.
    Moreover, 
    \[
        \big\lvert\,d(\exp_x\xi,\exp_x\eta)-\lvert\xi-\eta\rvert\,\big\rvert \le C_\ast\max(\lvert\xi\rvert,\lvert\eta\rvert)^2\lvert\xi-\eta\rvert 
    \] 
    for \(\lvert \xi \rvert, \lvert \eta\rvert \le r_1\). 
\end{lem}

\begin{proof}[Proof of Lemma~\ref{lem:exp.estimates}]
    Fix \(x \in M\).
    Our argument is based on the use of normal charts around \(x\).
    Write \( R\coloneqq \max(\lvert\xi\rvert,\lvert\eta\rvert)\le r_1 \), let \( g^x\coloneqq (\exp_x)^\ast g \) be the pullback of the metric to \( T_xM \), and identify \( T_p(T_xM)\cong T_xM \) canonically, so that the differential \( D(\exp_x)_p \) is a linear map \( T_xM\to T_{\exp_xp}M \).

    Consider the smooth function \(G(x,p,w)\coloneqq \big\lvert D(\exp_x)_p\,w\big\rvert^2\) on \(TM \oplus TM\).
    For fixed \( (x,w) \) the function \( p\mapsto G(x,p,w) \) satisfies
    \[
        G(x,0,w)=\lvert w\rvert^2,\qquad \partial_pG(x,0,w)=0 ,
    \]
    the first because \( D(\exp_x)_0=\Id \), the second because the Christoffel symbols vanish at the centre of a geodesic normal chart (remark both derivatives here are taken in the fibre \( T_xM \), which is a vector space, so they are intrinsic; no choice of frame is involved).  
    Taylor's theorem along the ray \( s\mapsto(x,sp,w) \) therefore gives
    \[
        \big\lvert G(x,p,w)-\lvert w\rvert^2\big\rvert
        \le\tfrac12\left(\sup_{\mathcal K}\big\lVert\partial^2_pG\big\rVert\right)\lvert p\rvert^2 ,
    \]
    where \(\mathcal K\coloneqq \{(x,p,w):\lvert p\rvert\le4r_1,\ \lvert w\rvert\le1\}\) and the supremum is finite because \( \partial^2_pG \) is continuous and \( \mathcal K\subset TM\oplus TM \) is compact.  
    Calling \( c_M \) that supremum and using homogeneity in \( w \), we get \( \big\lvert G(x,p,w)-\lvert w\rvert^2\big\rvert\le c_M\lvert p\rvert^2\lvert w\rvert^2 \) for all \( w \).
    Factoring the difference of squares gives us, for \(\lvert p\rvert\le4r_1\)
    \begin{equation}\label{eq:infinitesimal.comparison}
    \begin{split}
        \Big\lvert\;\big\lvert D(\exp_x)_p\,w\big\rvert-\lvert w\rvert\;\Big\rvert &= \frac{\big\lvert G(x,p,w)-\lvert w\rvert^2\big\rvert}{\big\lvert D(\exp_x)_p\,w\big\rvert+\lvert w\rvert}\\
        &\le \frac{\big\lvert G(x,p,w)-\lvert w\rvert^2\big\rvert}{\lvert w\rvert} \\ &\le\ c_M\lvert p\rvert^2\lvert w\rvert. 
    \end{split}
    \end{equation}
    Equivalently: on \( B(0,4r_1)\subset T_xM \) the metric \( g^x \) is within a factor \( 1\pm c_M\lvert p\rvert^2 \) of the Euclidean one. 
    This translates to distances on \(T_xM\) in the following way.
    For \( 0<r\le3r_1 \) let \( d^x_r \) denote the length distance of \( \left(\overline B(0,r),g^x\right) \).
    We claim that, for \( \xi,\eta\in\overline B(0,r) \),
    \begin{equation}\label{eq:ball.comparison}
        \left(1-c_Mr^2\right)\lvert\xi-\eta\rvert
        \ \le\ d^x_r(\xi,\eta)
        \ \le\ \left(1+c_Mr^2\right)\lvert\xi-\eta\rvert .
    \end{equation}
    For the upper bound, the Euclidean segment from \( \xi \) to \( \eta \) stays in \( \overline B(0,r) \) by convexity, and \eqref{eq:infinitesimal.comparison} bounds its \( g^x \)-length by \( (1+c_Mr^2)\lvert\xi-\eta\rvert \). 
    For the lower bound, we may assume \( c_Mr^2<1 \), the claim being vacuous otherwise.
    Then every Lipschitz curve in \( \overline B(0,r) \) joining \( \xi \) to \( \eta \) has, by \eqref{eq:infinitesimal.comparison}, \( g^x \)-length at least \( (1-c_Mr^2) \) times its Euclidean length, hence at least \( (1-c_Mr^2)\lvert\xi-\eta\rvert \).
    The distance \(d^x_r(\xi,\eta)\) is the infimum of all such lengths, and the claim follows.  
    Note that no curve here is required to be geodesic.
 
    Finally, we translate the bounds above to distances in \(M\).
    Since \( 4r_1<r_0 \), the map \( \exp_x \) restricts to an isometry from \( \left(B(0,4r_1),g^x\right) \) onto \( \left(B(x,4r_1),g\right) \), and \( \lvert\exp_x^{-1}y\rvert=d(x,y) \) there by the Gauss lemma.
    Let \( \gamma:[0,1]\to M \) be a minimising geodesic from \( \exp_x\xi \) to \( \exp_x\eta \).    
    From the triangle inequality~\eqref{eq:triangle.inequality}, \( L_g(\gamma)=d(\exp_x\xi,\exp_x\eta)\le\lvert\xi\rvert+\lvert\eta\rvert\le2r_1 \), so for every \( s \)
    \[
        d\left(x,\gamma(s)\right)\le d(x,\exp_x\xi)+L_g(\gamma)\le3r_1<4r_1 ;
    \]
    hence \( \widetilde\gamma\coloneqq \exp_x^{-1}\circ\,\gamma \) is a well-defined curve in \( \overline B(0,3r_1) \) joining \( \xi \) to \( \eta \).  
    The upper bound in \eqref{eq:ball.comparison}, applied with \( r=R \), gives
    \[
        \lvert\widetilde\gamma(s)\rvert=d\left(x,\gamma(s)\right)
        \le\lvert\xi\rvert+\left(1+c_MR^2\right)\lvert\xi-\eta\rvert
        \le\left(3+2c_MR^2\right)R\le C_2R,
    \]
    where \( C_2\coloneqq 3+2c_Mr_1^2 \).
    Put \( r\coloneqq \min\{C_2R,\,3r_1\} \).  
    Then \( L_g(\gamma) = L_{g^x}(\widetilde\gamma) \ge d^x_r(\xi,\eta) \), so \eqref{eq:ball.comparison} and \( r\le C_2R \) give
    \[
        d(\exp_x\xi,\exp_x\eta)\ \ge\ d^x_r(\xi,\eta)
        \ \ge\ \left(1-c_MC_2^2R^2\right)\lvert\xi-\eta\rvert .
    \]
    Putting this together with the upper bound of \eqref{eq:ball.comparison} evaluated at \( r=R \), we obtain
    \[
        \left(1-c_MC_2^2R^2\right)\lvert\xi-\eta\rvert \le d(\exp_x\xi,\exp_x\eta)\le (1+c_MR^2)\lvert\xi-\eta\rvert,
    \]  
    and by setting \( C_\ast\coloneqq c_MC_2^2=c_M(3+2c_Mr_1^2)^2 \) in the above inequalities we get \( \big\lvert\,d(\exp_x\xi,\exp_x\eta)-\lvert\xi-\eta\rvert\,\big\rvert \le C_\ast\max(\lvert\xi\rvert,\lvert\eta\rvert)^2\lvert\xi-\eta\rvert \) for \(\lvert \xi \rvert, \lvert \eta\rvert \le r_1\), as we wanted.
    Finally, take \(\Lambda\coloneqq 1+C_\ast r_1^2\) to obtain \(\exp_x: B(0, r_1) \to M\) \(\Lambda\)-Lipschitz for any \(x\in M\).
\end{proof}

\end{document}
